\documentclass[11pt,a4paper]{article}
\usepackage{amsmath,amssymb,amsthm}
\usepackage[margin=2.5cm]{geometry}
\usepackage{hyperref}

\newtheorem{definition}{Definition}[section]
\newtheorem{theorem}[definition]{Theorem}
\newtheorem{proposition}[definition]{Proposition}
\newtheorem{lemma}[definition]{Lemma}
\newtheorem{corollary}[definition]{Corollary}
\newtheorem{conjecture}[definition]{Conjecture}
\newtheorem{remark}[definition]{Remark}

\title{Hyperseed Formalization:\\
  Tensor Logic for Bridging Neural and Symbolic AI}
\author{ProtomegaTron (automated formalization)}
\date{2026-09-18}

\begin{document}
\maketitle

\begin{abstract}
We formalize the intellectual structure of Ben Goertzel's Substack article
``Tensor Logic for Bridging Neural and Symbolic AI'' (2025-12-16) within
the Hyperseed ontology. The article describes how Pedro Domingos's tensor
logic can serve as a mathematical lingua franca bridging symbolic and
neural AI within the Hyperon project. We model tensor logic as the
transition function between the neural and symbolic fibers in the
cognitive architecture bundle, identify curvature in the
representation-translation space, and show how tensor contractions
implement semantic transport.
\end{abstract}

\tableofcontents

%% =================================================================
\section{The neural--symbolic fiber bundle}
\label{sec:bundle}
%% =================================================================

\begin{definition}[Cognitive representation bundle --- claims 1--2]
\label{def:rep-bundle}
Let $\pi : E \to B$ be a fiber bundle where:
\begin{itemize}
\item $B$ = the space of cognitive contents (facts, rules, patterns,
  inferences),
\item $F_{\text{sym}}$ = the symbolic representation fiber (graphs, trees,
  logical formulas, pattern-matching structures),
\item $F_{\text{neur}}$ = the neural representation fiber (dense
  matrices, tensors, continuous activations on GPUs),
\item $E_b = F_{\text{sym}} \times F_{\text{neur}}$ for each cognitive
  content $b$.
\end{itemize}
\end{definition}

\begin{proposition}[Integration bottleneck --- claim 2]
\label{prop:bottleneck}
Current hybrid systems define ad hoc transition functions
$\tau : F_{\text{sym}} \to F_{\text{neur}}$ that are:
\begin{enumerate}
\item \textbf{Lossy}: symbolic structure is discarded during embedding.
\item \textbf{Serialized}: the translation is a sequential bottleneck
  between parallel GPU computation and sequential symbolic pattern
  matching.
\item \textbf{Non-compositional}: $\tau(A \circ B) \neq \tau(A) \circ \tau(B)$
  in general --- composing symbolic operations and then translating
  gives different results from translating and then composing neural
  operations.
\end{enumerate}
The cocycle condition fails: the transition functions don't compose
properly, which is why hybrid systems ``leave performance on the table.''
\end{proposition}

%% =================================================================
\section{Tensor logic as exact transition function}
\label{sec:tensor-logic}
%% =================================================================

\begin{definition}[Tensor logic embedding --- claims 4--8]
\label{def:tensor-embed}
Tensor logic defines a transition function
$\tau_{\text{TL}} : F_{\text{sym}} \to F_{\text{neur}}$ by:
\begin{enumerate}
\item \textbf{Facts $\to$ sparse tensors}: A logical database
  $\mathcal{D} = \{R(a_1, \ldots, a_k)\}$ maps to a sparse Boolean
  tensor $\mathbf{T}_R$ where $\mathbf{T}_R[a_1, \ldots, a_k] = 1$
  iff $R(a_1, \ldots, a_k) \in \mathcal{D}$.
\item \textbf{Rules $\to$ tensor contractions}: A logical rule
  $H(\vec{x}) \leftarrow B_1(\vec{y}_1), \ldots, B_m(\vec{y}_m)$
  maps to:
  \[
    \mathbf{T}_H[\vec{x}] = H\!\left(
      \sum_{\vec{z}} \prod_{i=1}^m \mathbf{T}_{B_i}[\vec{y}_i]
    \right)
  \]
  where $\vec{z}$ are the contracted (existentially quantified)
  variables and $H$ is a threshold function.
\item \textbf{Inference $\to$ einsum}: The contraction is an Einstein
  summation, directly executable on GPU tensor cores.
\end{enumerate}
\end{definition}

\begin{theorem}[Cocycle satisfaction for first-order logic --- claim 10]
\label{thm:cocycle}
For first-order Datalog (no function symbols, no negation), the tensor
logic transition function satisfies the cocycle condition:
\[
  \tau_{\text{TL}}(R_1 \bowtie R_2) = \tau_{\text{TL}}(R_1) \cdot \tau_{\text{TL}}(R_2)
\]
where $\bowtie$ is the logical join and $\cdot$ is tensor contraction.
Symbolic composition followed by translation equals translation followed
by neural composition. The fiber bundle is locally trivial in this
regime --- the transition is exact, not approximate.
\end{theorem}

\begin{corollary}[Seamless computational graph --- claim 8]
Because $\tau_{\text{TL}}$ satisfies the cocycle condition, neural
computations and logical inference can coexist in the same computational
graph without translation bottlenecks. Gradients can flow from neural
loss functions through logical inference steps and back.
\end{corollary}

%% =================================================================
\section{Curvature in representation space}
\label{sec:curvature}
%% =================================================================

\begin{definition}[Representation curvature --- claim 15]
\label{def:curvature}
Define the \emph{representation curvature} $\kappa(b)$ for a cognitive
content $b$ as the degree to which the tensor logic transition function
deviates from exactness:
\[
  \kappa(b) = \|\tau_{\text{TL}}(\text{sym}(b)) - \text{neur}(b)\|
\]
where $\text{sym}(b)$ and $\text{neur}(b)$ are the ideal symbolic and
neural representations.
\end{definition}

\begin{proposition}[Flat and curved regions --- claims 4, 15]
\label{prop:flat-curved}
The representation space has identifiable regions of different curvature:
\begin{enumerate}
\item \textbf{Flat} ($\kappa \approx 0$): First-order logic, Datalog,
  ground facts, simple joins. The tensor logic embedding is exact.
\item \textbf{Low curvature}: Probabilistic/fuzzy extensions (PLN-style
  truth values). Replace Boolean tensors with $[0,1]$-valued tensors;
  the contraction semantics extends naturally.
\item \textbf{High curvature}: Higher-order logic, meta-reasoning,
  self-referential inference. The tensor representation requires
  approximation or architectural extension (e.g., tensor networks,
  attention mechanisms).
\end{enumerate}
\end{proposition}

\begin{remark}[Research frontier as curvature boundary]
The research frontier of tensor logic is precisely the boundary between
low and high curvature regions. The flat region is well-understood and
GPU-ready. The high-curvature region is where the most interesting
cognitive capabilities live (meta-reasoning, analogy, self-reflection)
and where new mathematical machinery is needed.
\end{remark}

%% =================================================================
\section{PLN as probabilistic fiber extension}
\label{sec:pln}
%% =================================================================

\begin{definition}[PLN tensor extension --- claim 14]
\label{def:pln-tensor}
Extend the Boolean tensor framework to PLN by replacing
$\{0, 1\}$-valued tensors with tensors valued in the PLN truth-value
space $\mathcal{T} = [0,1] \times [0,1]$ (frequency $\times$ confidence):
\[
  \mathbf{T}_R^{\text{PLN}}[a_1, \ldots, a_k] = (f, c) \in \mathcal{T}
\]
Tensor contractions become probabilistic inference rules:
\[
  (f_H, c_H) = \text{PLN-rule}\!\left(
    \bigoplus_{\vec{z}} \bigotimes_{i=1}^m (f_{B_i}, c_{B_i})
  \right)
\]
where $\bigoplus$ and $\bigotimes$ are the PLN combination operators.
\end{definition}

\begin{proposition}[Unified inference pipeline]
\label{prop:unified}
The PLN tensor extension enables a unified inference pipeline where:
\begin{enumerate}
\item Neural pattern recognition produces $(f, c)$ estimates.
\item Logical rules compose these estimates via tensor contractions.
\item Uncertainty propagates naturally through the contraction.
\item The entire pipeline runs on GPU tensor cores.
\end{enumerate}
This is the concrete mechanism for cognitive synergy between neural
perception and symbolic reasoning.
\end{proposition}

%% =================================================================
\section{Hyperon architecture mapping}
\label{sec:hyperon}
%% =================================================================

\begin{definition}[Hyperon component mapping --- claims 12--13]
\label{def:hyperon-map}
The Hyperon architecture maps onto the tensor logic framework:
\begin{itemize}
\item \textbf{MORK} (metagraph database) $\leftrightarrow$ sparse tensor
  storage. MORK exports logical facts as sparse tensors for GPU inference
  and imports results back into the metagraph.
\item \textbf{PeTTa/MM2} (MeTTa interpreters) $\leftrightarrow$ tensor
  contraction compilers. MeTTa pattern-matching rules compile to einsum
  specifications.
\item \textbf{Neural modules} $\leftrightarrow$ dense tensor operations.
  Standard deep learning components run on the same GPU.
\item \textbf{Tensor logic} $\leftrightarrow$ transition function
  $\tau_{\text{TL}}$ connecting all three.
\end{itemize}
\end{definition}

%% =================================================================
\section{d-calculus perspective}
\label{sec:d-calc}
%% =================================================================

\begin{remark}[Tensor contractions as semantic transport]
In the d-calculus framework, a tensor contraction implements semantic
transport: it moves information from the symbolic domain to the neural
domain (or vice versa) while preserving inferential structure. The
contraction indices are the transport coordinates --- they specify
which dimensions of meaning are preserved and which are summed over
(projected out). This is precisely the fiber-bundle transport that
Hyperseed requires for cross-paradigm inference.
\end{remark}

\begin{remark}[GPU democratization as base-space expansion]
If logical reasoning can run on GPUs (claim 21), the base space of
available computational substrates for symbolic AI expands dramatically.
The massive hardware investment in tensor processors --- originally
built for the neural fiber only --- becomes available to the symbolic
fiber. This is a base-space topology change at the hardware level,
analogous to the organizational topology change described in the
OpenBGI article.
\end{remark}

\begin{remark}[Unification as shared fiber structure]
The article's deepest claim (claim 19) is that the apparent gulf between
symbolic and neural AI is an artifact of implementation history, not
fundamental mathematics. In Hyperseed terms: the neural and symbolic
fibers are not independent --- they share a common sub-fiber (tensor
algebra) that has been obscured by historical accident. Tensor logic
makes this shared structure explicit, enabling genuine integration
rather than ad hoc bridging.
\end{remark}

\end{document}
